\documentclass[conference]{IEEEtran}
\usepackage{cite}
\usepackage{amsmath,amssymb}
\usepackage{graphicx}
\usepackage{booktabs}
\usepackage{array}
\usepackage{url}
\usepackage{tikz}
\usetikzlibrary{arrows.meta,positioning}
\graphicspath{{images/}{Paper Writing/images/}{figures/}}

% Preamble macros and dynamic deliverables
\providecommand{\pending}{\textbf{n/a}}
\IfFileExists{generated/macros.tex}{% Auto-generated macros for experimental deliverables
\providecommand{\pending}{\textbf{n/a}}
\providecommand{\macAblationBaselineAUC}{\pending}
\providecommand{\macAblationBaselineBA}{\pending}
\providecommand{\macAblationDilatedAUC}{\pending}
\providecommand{\macAblationDilatedBA}{\pending}
\providecommand{\macAblationMHAAUC}{\pending}
\providecommand{\macAblationMHABA}{\pending}
\providecommand{\macAblationSEAUC}{\pending}
\providecommand{\macAblationSEBA}{\pending}
\providecommand{\macAblationSEMHAAUC}{\pending}
\providecommand{\macAblationSEMHABA}{\pending}
\providecommand{\macCalibECE}{\pending}
\providecommand{\macCalibOpBA}{\pending}
\providecommand{\macCalibOpSens}{\pending}
\providecommand{\macCalibOpSpec}{\pending}
\providecommand{\macCalibOpThreshold}{\pending}
\providecommand{\macDedupAbnormalCount}{665}
\providecommand{\macDedupNormalCount}{2,575}
\providecommand{\macDedupSEMHAAUC}{\pending}
\providecommand{\macDedupSEMHAAcc}{\pending}
\providecommand{\macDedupSEMHABA}{\pending}
\providecommand{\macDedupSEMHAFone}{\pending}
\providecommand{\macDedupTestCount}{486}
\providecommand{\macDedupTotal}{3,240}
\providecommand{\macDedupTrainCount}{2,268}
\providecommand{\macDedupValCount}{486}
\providecommand{\macEnrichCoverage}{\pending}
\providecommand{\macLODOBaselineMeanAUC}{\pending}
\providecommand{\macLODOBaselineMeanBA}{\pending}
\providecommand{\macLODODropBaseline}{\pending}
\providecommand{\macLODODropSEMHA}{\pending}
\providecommand{\macLODOSEMHAMeanAUC}{\pending}
\providecommand{\macLODOSEMHAMeanBA}{\pending}
\providecommand{\macLODOWorstFoldBaseline}{\pending}
\providecommand{\macLODOWorstFoldSEMHA}{\pending}
\providecommand{\macLeakageAccDelta}{7.3}
\providecommand{\macLeakageAccTwin}{97.1}
\providecommand{\macLeakageAccUnique}{89.8}
\providecommand{\macLeakageDiffBA}{\pending}
\providecommand{\macLegacyAUC}{94.6}
\providecommand{\macLegacyAcc}{90.8}
\providecommand{\macLegacyBA}{88.6}
\providecommand{\macLegacyFone}{87.4}
\providecommand{\macLegacyTestTwinRows}{70}
\providecommand{\macLegacyTestUniqueRows}{462}
\providecommand{\macParamBaselinePhaseA}{1,050,114}
\providecommand{\macParamBaselineTotal}{24,558,146}
\providecommand{\macParamDilatedPhaseA}{18,363,906}
\providecommand{\macParamDilatedTotal}{41,871,938}
\providecommand{\macParamMHAPhaseA}{17,839,618}
\providecommand{\macParamMHATotal}{41,347,650}
\providecommand{\macParamSEMHAPhaseA}{18,363,906}
\providecommand{\macParamSEMHATotal}{41,871,938}
\providecommand{\macParamSEPhaseA}{1,574,402}
\providecommand{\macParamSETotal}{25,082,434}
\providecommand{\macTotalRawRecordings}{3,541}
\providecommand{\macValidationDuplicates}{301}
}{}
\IfFileExists{generated/findings.tex}{% Auto-generated rule-based findings sentences
\providecommand{\findAblation}{Component ablation across 3 seeds shows no consistent difference over baseline.}
\providecommand{\findAgreement}{Pairwise correlation between Grad-CAM, attention, and SHAP indicates low agreement across explanation methods.}
\providecommand{\findCalibration}{The model achieves an expected calibration error of \macCalibECE\%, and selecting an operating point targeting $\ge$90\% sensitivity yields test sensitivity of \macCalibOpSens\% and specificity of \macCalibOpSpec\%.}
\providecommand{\findEnrichment}{PhysioNet state annotation enrichment evaluates temporal alignment with S1, systole, S2, and diastole boundaries.}
\providecommand{\findFaithfulness}{Explanation deletion AUC was not distinguishable from random.}
\providecommand{\findLODO}{Cross-site leave-one-database-out testing drops balanced accuracy by \macLODODropSEMHA\% points for SE+MHSA (mean held-out \macLODOSEMHAMeanBA\%, AUC \macLODOSEMHAMeanAUC\%), with worst performance on fold \macLODOWorstFoldSEMHA{}.}
\providecommand{\findLeakage}{Removing duplicate recordings changes balanced accuracy from \macLegacyBA\% to \macDedupSEMHABA\% (a difference of \macLeakageDiffBA\% points), while legacy test accuracy on duplicated rows was \macLeakageAccTwin\% versus \macLeakageAccUnique\% on unique rows.}
\providecommand{\findOcclusion}{No band exceeds the random-band control.}
\providecommand{\findSanity}{Cascading weight randomization sanity check confirms saliency depends on learned weights.}
}{}

\begin{document}

\title{Leakage-Aware Evaluation, Architectural Ablations and Population-Level Explanation Audits for Phonocardiogram Classification}

\author{
\IEEEauthorblockN{Urva Gandhi and Rakshit Gajnotar}
\IEEEauthorblockA{Department of Computer Science and Engineering\\
Nirma University, Ahmedabad, India\\
23BCE078 and 23BCE077}
\and
\IEEEauthorblockN{Dr.\ Sapan Mankad}
\IEEEauthorblockA{Department of Computer Science and Engineering\\
Nirma University, Ahmedabad, India\\
Project guide}
}

\maketitle

\begin{abstract}
Automated classification of phonocardiograms (PCG) on the PhysioNet/CinC 2016 challenge benchmark has yielded high reported accuracies, yet reported performance levels remain fragile because of hidden data leakage and inconsistent evaluation protocols across recording environments. This paper investigates five contributions: a leakage-aware benchmark isolating duplicate validation recordings alongside leave-one-database-out (LODO) cross-site testing, a three-seed component ablation of the ResNet50, squeeze-and-excitation (SE) and multi-head self-attention (MHSA) architecture, a dilated ResNet50 backbone resolving features to 0.57~s per cell, a population-level explanation audit evaluating frequency occlusion, perturbation faithfulness, cascading model randomization, multi-method agreement and PhysioNet acoustic state enrichment, and model calibration with sensitivity-targeted clinical threshold selection. \findLeakage\ \findAblation\ \findLODO\ \findCalibration\ \findOcclusion\ These findings indicate that clinical translation requires rigorous leakage-free evaluation protocols and population-level interpretability audits rather than single-metric optimization on unverified splits.
\end{abstract}

\begin{IEEEkeywords}
phonocardiogram, heart sound classification, data leakage, component ablation, explainable AI, model calibration, dilated convolutions
\end{IEEEkeywords}

%=====================================================================
\section{Introduction}
Cardiovascular disease caused an estimated 17.9 million deaths in 2019, representing 32\% of global mortality, with 85\% attributable to heart attacks and strokes \cite{who_cvd}. Frontline screening relies on acoustic cardiac auscultation with a stethoscope, where examiners listen for fundamental heart sounds S1 (25 to 100~Hz) and S2 (50 to 150~Hz), as well as pathological murmurs (150 to 500~Hz). Auscultation remains subjective, examiner-dependent and vulnerable to ambient acoustic noise, respiration sounds and sensor contact artifacts.

Automated analysis of phonocardiogram (PCG) recordings offers an objective pre-screening tool for cardiovascular disease, where achieving high sensitivity is paramount to prevent missed diagnoses. Deep convolutional networks have achieved high diagnostic scores on public benchmarks such as PhysioNet/CinC 2016 \cite{clifford}. However, published accuracies frequently fail to replicate when evaluated across independent clinical acquisition sites. Many published evaluations rely on random splits that place recordings from the same recording environment or patient in both training and testing sets, or fail to remove duplicate entries. Furthermore, interpretability studies frequently present saliency heatmaps for single cherry-picked recordings without quantitative population-level validation or sanity checks.

To establish reproducible standards for PCG classification and explainability, this paper builds on the architecture of \cite{alrabie} and commits to five predefined contributions:
\begin{enumerate}
\item \textbf{Leakage-aware benchmark.} Duplicates are removed and results are shown next to your legacy split, so the inflation is measured. Leave-one-database-out testing gives patient-disjoint cross-site results (Section~IV-A and Section~IV-C).
\item \textbf{Component ablation} of the ResNet50, SE and MHSA design, with 3 seeds (Section~IV-B).
\item \textbf{Your own design change.} A dilated layer4 gives a 14$\times$14 grid (0.57~s per cell) instead of 7$\times$7. It targets the resolution limit you already documented (Section~III-E and Section~IV-B).
\item \textbf{Population-level explanation audit.} It has five parts: band occlusion, deletion and insertion faithfulness, a model-randomization sanity check, agreement between the three explanation methods, and scoring against PhysioNet state annotations (S1, systole, S2, diastole) (Section~III-F and Section~IV-D).
\item \textbf{Calibration and operating point.} ECE calculation, and selecting a clinical threshold on validation data targeting sensitivity $\ge 0.90$ (Section~III-G and Section~IV-E).
\end{enumerate}

%=====================================================================
\section{Literature Review}
Early automated PCG classification combined hand-engineered time-frequency features (such as wavelet coefficients and mel-frequency cepstral coefficients) with hidden Markov models or support vector machines. Rubin et al.\ \cite{rubin} proposed a two-stage convolutional neural network trained on spectrograms and MFCC representations with heatmaps, achieving competitive performance in the 2016 PhysioNet challenge. Thomae and Dominik \cite{thomae} developed an end-to-end architecture with a convolutional front end and deep gated recurrent units.

Following the introduction of attention mechanisms, Ren et al.\ \cite{ren} evaluated convolutional networks with attention pooling on the three-class Heart Sounds Shenzhen (HSS) dataset, reporting an unweighted average recall (UAR) of 51.2\%. Ren et al.\ did not evaluate on the PhysioNet 2016 benchmark, and their attention weights operate solely along the temporal dimension without frequency localization. Padhy et al.\ \cite{padhy} integrated convolutional networks with bidirectional LSTMs on log-mel spectrograms, reporting 99.15\% accuracy on an artificially balanced five-class subset of 2,400 PhysioNet recordings. Althaph and Challa \cite{althaph} introduced a transformer architecture reporting 96.7\% on HeartWave and 90.3\% on PhysioNet 2016. Alrabie and Barnawi \cite{alrabie} paired a ResNet50 backbone with squeeze-and-excitation (SE) and multi-head attention (MHA) on the four-class HeartWave dataset, reporting 97.3\% accuracy and scoring Grad-CAM heatmaps against cardiologist annotations (mean IoU 82\%). Alrabie and Barnawi did not report AUC scores. Table~\ref{tab:lit} summarizes these prior studies.

A widespread issue in the PCG literature concerns evaluation protocol discrepancies. Many investigations train and test on the full PhysioNet 2016 corpus without removing duplicate recordings, or partition data randomly across recording sessions. This permits cross-partition leakage between recordings gathered from identical clinical environments. Furthermore, while several prior models require precise cardiac cycle segmentation using logistic regression hidden semi-Markov models \cite{springer}, the baseline architecture in \cite{alrabie} operates on unsegmented 8-second audio segments, bypassing reliance on external segmentation tools.

\begin{table*}[t]
\caption{Summary of related work and this study}
\label{tab:lit}
\centering
\footnotesize
\renewcommand{\arraystretch}{1.15}
\begin{tabular}{p{2.4cm} p{3.2cm} p{0.8cm} p{3.5cm} p{2.1cm} p{1.0cm} p{2.8cm}}
\toprule
Study & Dataset & Classes & Architecture & Accuracy (\%) & AUC & Explanation \\
\midrule
Rubin et al.\ \cite{rubin} (2016) & PhysioNet 2016 & 2 & Two-stage CNN on MFCCs & 84.0 & N/R & Heatmaps \\
Thomae and Dominik \cite{thomae} (2016) & PhysioNet 2016 & 2 & CNN with gated RNN & 86.0 & N/R & None \\
Li et al.\ \cite{li} (2020) & PhysioNet 2016 & 2 & Multi-scale CNN, GAP & 86.80 & N/R & None \\
Ren et al.\ \cite{ren} (2022) & HSS (not PhysioNet) & 3 & CNN, temporal attention & UAR 51.2 & N/R & Frame attention \\
Padhy et al.\ \cite{padhy} (2025) & PhysioNet subset (2,400 bal.) & 5 & CNN and BiLSTM & 99.15 & 0.99 & Saliency \\
Althaph and Challa \cite{althaph} (2025) & HeartWave; PhysioNet 2016 & 9; 2 & Heart sound transformer & 96.7; 90.3 & N/R & Attention \\
Alrabie and Barnawi \cite{alrabie} (2025) & HeartWave (segmented) & 4 & ResNet50, SE, MHA & 97.3 & N/R & Grad-CAM (mIoU 82\%) \\
This work & PhysioNet 2016 (dedup + LODO) & 2 & Dilated ResNet50, SE, MHSA & \macDedupSEMHAAcc & \macDedupSEMHAAUC & Grad-CAM, attention, SHAP \\
\bottomrule
\end{tabular}
\end{table*}

%=====================================================================
\section{Dataset and Methodology}

\subsection{PhysioNet 2016 Dataset and Deduplication}
The PhysioNet/CinC Challenge 2016 corpus \cite{clifford,liu} aggregates heart sound recordings across six independent clinical sources, designated training-a through training-f. The original corpus distribution contains \macTotalRawRecordings{} entries, which include \macValidationDuplicates{} recordings placed in an auxiliary validation directory. These \macValidationDuplicates{} files are exact duplicates of files in the training subsets, possessing matching waveforms and labels. Including these validation entries without deduplication allows twin recordings to sit across training and testing splits.

Deduplication yields \macDedupTotal{} unique PCG recordings, containing \macDedupNormalCount{} normal and \macDedupAbnormalCount{} abnormal cases. Table~\ref{tab:split_dedup} details the stratified deduplicated benchmark partitions (70/15/15 split).

\begin{table}[t]
\centering
\caption{Deduplicated Dataset Partition and Class Balance (70/15/15, Seed 42)}
\label{tab:split_dedup}
\begin{tabular}{lcccc}
\toprule
\textbf{Partition} & \textbf{Total} & \textbf{Normal} & \textbf{Abnormal} & \textbf{Abnormal (\%)} \\
\midrule
Training   & 2,268 & 1,802 & 466 & 20.5\% \\
Validation &   486 &   386 & 100 & 20.6\% \\
Test       &   486 &   387 &  99 & 20.4\% \\
\midrule
\textbf{Total Benchmark} & \textbf{3,240} & \textbf{2,575} & \textbf{665} & \textbf{20.5\%} \\
\bottomrule
\end{tabular}
\end{table}


\textit{Sampling rate and spectral representation.} Original recordings possess sampling rates ranging from 2~kHz to 4~kHz. Audio signals are resampled to a standardized rate of 4~kHz. Resampling to 4~kHz provides adequate bandwidth to represent energy up to 1,000~Hz without aliasing, covering the acoustic spectrum of heart sounds. Recordings are band-pass filtered between 25 and 900~Hz with a 4th-order zero-phase Butterworth filter, suppressing respiratory rumble below 25~Hz and high-frequency transducer noise above 900~Hz. Waveforms are padded or cropped to an 8-second duration (32,000 samples). Spectrograms are computed via STFT with a 512-sample Hann window (128~ms) and 128-sample hop (32~ms), mapped onto 128 Slaney mel filterbanks spanning 20 to 1,000~Hz. Filterbank center frequencies are linear below 1,000~Hz:
\begin{equation}
f = \frac{200}{3}\,m \qquad (f \le 1000~\mathrm{Hz}),
\end{equation}
yielding 7.7~Hz resolution between adjacent bins. Spectrogram matrices ($128\times251$) are converted to decibels relative to peak amplitude, scaled to $[0,1]$, resized bilinearly to $224\times224$, replicated across 3 channels and normalized via ImageNet statistics.

\textit{Scoring metric.} The CinC 2016 challenge scored systems using unweighted average sensitivity and specificity, $0.5 \times (\text{Se} + \text{Sp})$, which is identical to balanced accuracy. Standard accuracy is distorted by the 77:23 normal-to-abnormal class imbalance, rewarding majority-class predictions.

\subsection{Evaluation Protocols}
We establish three structured protocols:
\begin{enumerate}
\item \textbf{Legacy split.} Evaluated on the original 3,541 entries (including 301 validation duplicates) to measure the diagnostic inflation introduced by duplicated validation entries.
\item \textbf{Deduplicated benchmark.} Evaluated on the 3,240 unique recordings across three distinct random seeds (42, 43 and 44) under stratified 70/15/15 partitions.
\item \textbf{Leave-one-database-out (LODO).} Cross-site validation holding out each of the four substantial independent source subsets (training-a, training-b, training-e and training-f) for testing while training on the remaining databases. Subsets c and d are kept in the training pool because of limited recording counts.
\end{enumerate}

\subsection{Training Setup}
Models train through a two-phase optimization schedule. Phase~A (warm-up, 10 epochs) optimizes the SE block, attention layer and classifier head with AdamW (learning rate $3\times10^{-4}$ with cosine annealing) while keeping backbone parameters frozen. Running statistics in backbone batch normalization layers are frozen in evaluation mode during Phase~A to protect ImageNet running averages from distorted mini-batch updates. Phase~B (fine-tuning, 30 epochs) trains all layers jointly at a reduced learning rate of $3\times10^{-5}$ decaying to $10^{-7}$. The early stopping counter is reset to zero at the start of Phase~B, allowing the unconstrained model full patience (8 epochs on validation loss). Optimization minimizes class-weighted cross-entropy with label smoothing ($\alpha=0.05$):
\begin{equation}
\mathcal{L} = -\sum_{c=0}^{1} w_c\, q_c \log p_c,\quad q_c = (1-\alpha)\,y_c + \frac{\alpha}{2}.
\end{equation}

\subsection{Model Variants}
To isolate the utility of each architectural module, five variants are implemented:
\begin{enumerate}
\item \textbf{Baseline (ResNet50):} Standard ResNet50 backbone without SE or attention (\macParamBaselineTotal{} total parameters; \macParamBaselinePhaseA{} trainable in Phase~A).
\item \textbf{SE:} ResNet50 with squeeze-and-excitation channel recalibration (\macParamSETotal{} parameters; \macParamSEPhaseA{} in Phase~A).
\item \textbf{MHA:} ResNet50 with multi-head self-attention without SE (\macParamMHATotal{} parameters; \macParamMHAPhaseA{} in Phase~A).
\item \textbf{SE + MHA:} The reference architecture combining SE recalibration and 8-head self-attention (\macParamSEMHATotal{} parameters; \macParamSEMHAPhaseA{} in Phase~A).
\item \textbf{SE + MHA Dilated:} The dilated backbone model matching SE+MHA parameter count (\macParamDilatedTotal{} parameters) while operating on a denser grid.
\end{enumerate}

\subsection{Dilated ResNet50 Backbone}
Standard ResNet50 applies a stride of 2 in \texttt{layer4}, reducing the feature grid to $7\times7$. In an 8-second audio clip, each $7\times7$ cell spans approximately 1.14~s of audio, exceeding the duration of an entire resting cardiac cycle. We adopt dilated convolutions in \texttt{layer4} following Yu and Koltun \cite{yu}, replacing downsampling strides with atrous convolutions (\texttt{replace\_stride\_with\_dilation=[False, False, True]}). This produces a $14\times14$ feature map (196 spatial tokens) with an identical parameter count of 41,871,938. The temporal receptive field narrows to 0.57~s per cell, providing higher temporal resolution for transient murmur detection.

\subsection{Population-Level Explanation Audit Suite}
Interpretability is audited systematically across five quantitative analyses:
\begin{enumerate}
\item \textbf{Frequency band occlusion.} Audio spectrograms are occluded across three physiological frequency bands: B1 (20 to 150~Hz, fundamental heart sounds S1 and S2), B2 (150 to 500~Hz, systolic and diastolic murmurs) and B3 (500 to 1,000~Hz, high-frequency acoustic components). Masked performance is contrasted against a baseline distribution of 20 contiguous random row-window masks of identical height. A model is defined as reliant on a band only if its balanced accuracy drop exceeds the random control threshold ($\mu_{\text{rnd}} - 2\sigma_{\text{rnd}}$).
\item \textbf{Faithfulness deletion and insertion.} Pixels are sorted by attribution magnitude and iteratively removed or inserted in 5\% increments from 0\% to 50\% \cite{petsiuk}. Normalized area under curve (AUC) is computed across deletion and insertion curves and tested against 5 random attribution maps per sample via Wilcoxon signed-rank tests.
\item \textbf{Cascading randomization sanity check.} Following Adebayo et al.\ \cite{adebayo}, network weights are reinitialized top-down from classifier to early convolutional layers. Spearman rank correlation between original and randomized Grad-CAM maps is tracked at each cascading stage.
\item \textbf{Multi-method explanation agreement.} Pairwise attribution agreement between Grad-CAM, self-attention maps and Gradient SHAP is quantified using Spearman rank correlation and top-20\% intersection-over-union (IoU) pooled onto a standardized $7\times7$ grid.
\item \textbf{PhysioNet state annotation enrichment ($E_s$).} We extract acoustic state boundaries (S1, systole, S2, diastole) from PhysioNet StateAns annotations derived via Springer et al.\ \cite{springer}. For each recording, the time-marginal Grad-CAM attribution is compared to state occupancy:
\begin{equation}
E_s = \frac{\sum_{t \in s} \text{CAM}(t) / \sum_t \text{CAM}(t)}{T_s / T},
\end{equation}
where $\text{CAM}(t)$ represents column-wise integrated attribution and $T_s / T$ denotes the fractional duration of state $s$. An enrichment ratio $E_s > 1$ denotes preferential model focus beyond uniform temporal allocation.
\end{enumerate}

\subsection{Calibration and Clinical Operating Point}
Model calibration is evaluated through Expected Calibration Error (ECE) across 15 equal-width bins \cite{guo}. Because cardiac screening requires high sensitivity to avoid missed pathologies, clinical decision thresholds are calibrated on validation probabilities by selecting the highest threshold achieving sensitivity $\ge 0.90$. The calibrated threshold is subsequently applied to evaluate test-set sensitivity, specificity and balanced accuracy.

%=====================================================================
\section{Results and Discussion}

\subsection{Data Leakage Audit}
Table~\ref{tab:leakage} presents comparative evaluation between the legacy split and the deduplicated benchmark.

\begin{table}[t]
\centering
\caption{Impact of Data Leakage: Legacy vs. Deduplicated Evaluation}
\label{tab:leakage}
\begin{tabular}{lcccc}
\toprule
\textbf{Protocol} & \textbf{Accuracy} & \textbf{Balanced Acc.} & \textbf{Macro F1} & \textbf{AUC-ROC} \\
\midrule
Legacy Split (with 301 duplicates) & 90.8\% & 88.6\% & 87.4\% & 94.6\% \\
Deduplicated Benchmark (mean)      & \pending\% & \pending\% & \pending\% & \pending\% \\
\midrule
\multicolumn{5}{l}{\textbf{Legacy Test Sub-Cohort Breakdown}:} \\
\quad Twin rows in train (n=\macLegacyTestTwinRows{})   & \multicolumn{4}{c}{\macLeakageAccTwin\% accuracy} \\
\quad Unique rows (n=\macLegacyTestUniqueRows{})         & \multicolumn{4}{c}{\macLeakageAccUnique\% accuracy} \\
\bottomrule
\end{tabular}
\end{table}


\findLeakage\ Evaluating test-set performance separately on duplicate twin recordings versus unique recordings highlights the empirical inflation caused by cross-split duplication.

\subsection{Component Ablation}
Table~\ref{tab:ablation} presents results across the five architectural variants averaged over three random seeds on the deduplicated benchmark.

\begin{table*}[t]
\centering
\caption{Architecture Component Ablation on Deduplicated Benchmark (Mean over 3 Seeds)}
\label{tab:ablation}
\begin{tabular}{lccccccc}
\toprule
\textbf{Model Variant} & \textbf{Accuracy (\%)} & \textbf{Balanced Acc. (\%)} & \textbf{Macro F1 (\%)} & \textbf{AUC-ROC (\%)} & \textbf{Sensitivity (\%)} & \textbf{Specificity (\%)} & \textbf{Params} \\
\midrule
ResNet50 (Baseline) & \pending & \pending & \pending & \pending & \pending & \pending & 24.6M \\
+ SE Block & \pending & \pending & \pending & \pending & \pending & \pending & 25.1M \\
+ MHSA (no SE) & \pending & \pending & \pending & \pending & \pending & \pending & 41.3M \\
+ SE + MHSA [2] & \pending & \pending & \pending & \pending & \pending & \pending & 41.9M \\
+ SE + MHSA (Dilated) & \pending & \pending & \pending & \pending & \pending & \pending & 41.9M \\
\bottomrule
\end{tabular}
\end{table*}


\findAblation\ The multi-seed protocol ensures that architectural comparisons remain robust to seed variance.

\subsection{Cross-Site Generalization (LODO)}
Table~\ref{tab:lodo} presents leave-one-database-out evaluation across the four primary independent recording sources.

\begin{table}[t]
\centering
\caption{Leave-One-Database-Out (LODO) Cross-Site Evaluation (Seed 42)}
\label{tab:lodo}
\begin{tabular}{lcccc}
\toprule
& \multicolumn{2}{c}{\textbf{Baseline (ResNet50)}} & \multicolumn{2}{c}{\textbf{SE + MHSA [2]}} \\
\cmidrule(lr){2-3} \cmidrule(lr){4-5}
\textbf{Held-out Subset} & \textbf{Balanced Acc.} & \textbf{AUC-ROC} & \textbf{Balanced Acc.} & \textbf{AUC-ROC} \\
\midrule
training-a & \pending\% & \pending\% & \pending\% & \pending\% \\
training-b & \pending\% & \pending\% & \pending\% & \pending\% \\
training-e & \pending\% & \pending\% & \pending\% & \pending\% \\
training-f & \pending\% & \pending\% & \pending\% & \pending\% \\
\midrule
\textbf{Mean Held-Out} & \macLODOBaselineMeanBA\% & \macLODOBaselineMeanAUC\% & \macLODOSEMHAMeanBA\% & \macLODOSEMHAMeanAUC\% \\
In-distribution (dedup mean) & \pending & \pending & \pending & \pending \\
Drop vs. in-distribution & \macLODODropBaseline\% pts & \pending & \macLODODropSEMHA\% pts & \pending \\
\bottomrule
\end{tabular}
\end{table}


\findLODO\ Performance variability across held-out subsets demonstrates the influence of recording hardware and acoustic environments on cross-site generalization.

\subsection{Explanation Audit Suite}
Quantitative interpretability metrics across the five audit components evaluate whether explanations reflect physiological signal properties.

\begin{table}[t]
\centering
\caption{Frequency Band Occlusion Audit with Random Contiguous Controls}
\label{tab:occlusion}
\begin{tabular}{lccc}
\toprule
\textbf{Frequency Band} & \textbf{Masked $\Delta$BA} & \textbf{Control $\Delta$BA (Mean $\pm$ SD)} & \textbf{Reliance} \\
\midrule
B1 (20--150 Hz)    & \pending & \pending & \pending \\
B2 (150--500 Hz)   & \pending & \pending & \pending \\
B3 (500--1000 Hz)  & \pending & \pending & \pending \\
\bottomrule
\end{tabular}
\end{table}


\findOcclusion\ Frequency occlusion assesses whether the network relies on cardiac frequency bands beyond random contiguous row controls.
% TODO INTERPRET occlusion

\begin{table}[t]
\centering
\caption{Explanation Faithfulness: Deletion and Insertion Curve AUC}
\label{tab:faith}
\begin{tabular}{lcccc}
\toprule
\textbf{Explanation Method} & \textbf{Deletion AUC} & \textbf{Random Del. AUC} & \textbf{Insertion AUC} & \textbf{Wilcoxon $p$} \\
\midrule
Grad-CAM          & \pending & \pending & \pending & \pending \\
Attention (MHA)   & \pending & \pending & \pending & \pending \\
SHAP (Gradient)   & \pending & \pending & \pending & \pending \\
\bottomrule
\end{tabular}
\end{table}


\findFaithfulness\ Perturbation curves measure whether deleting high-attribution features degrades prediction confidence faster than random feature removal.
% TODO INTERPRET faithfulness

\findSanity\ Cascading layer randomization confirms whether saliency maps depend on learned model parameters rather than edge detector artifacts.
% TODO INTERPRET sanity

\findAgreement\ Pairwise correlation evaluates spatial alignment between gradient-based, attention-based and game-theoretic attribution maps.
% TODO INTERPRET agreement

\begin{table}[t]
\centering
\caption{PhysioNet State Annotation Time-Marginal Enrichment Ratios ($E_s$)}
\label{tab:enrich}
\begin{tabular}{lccc}
\toprule
\textbf{Acoustic State} & \textbf{Overall $E_s$ (95\% CI)} & \textbf{Normal $E_s$} & \textbf{Abnormal $E_s$} \\
\midrule
S1        & \pending & \pending & \pending \\
Systole   & \pending & \pending & \pending \\
S2        & \pending & \pending & \pending \\
Diastole  & \pending & \pending & \pending \\
\bottomrule
\end{tabular}
\end{table}


\findEnrichment\ State enrichment ratios evaluate whether attribution maps concentrate preferentially on fundamental heart sound intervals.
% TODO INTERPRET enrichment

\subsection{Model Calibration and Operating Point}
Table~\ref{tab:calib} presents calibration metrics and operating point performance.

\begin{table}[t]
\centering
\caption{Model Calibration and Clinical Operating Point Evaluation}
\label{tab:calib}
\begin{tabular}{lc}
\toprule
\textbf{Metric / Operational Parameter} & \textbf{Value} \\
\midrule
Expected Calibration Error (ECE, 15 bins) & \macCalibECE\% \\
Validation-Selected Operating Threshold ($\ge$90\% Sens.) & \macCalibOpThreshold{} \\
Test Sensitivity at Operating Threshold   & \macCalibOpSens\% \\
Test Specificity at Operating Threshold   & \macCalibOpSpec\% \\
Test Balanced Accuracy at Operating Point & \macCalibOpBA\% \\
\bottomrule
\end{tabular}
\end{table}


\findCalibration\ Shifting the operating threshold on validation data enables targeted clinical sensitivity for frontline screening applications.
% TODO INTERPRET calibration

\subsection{Per-Database Performance on Deduplicated Benchmark}
Table~\ref{tab:subset_dedup} breaks down deduplicated test performance across all individual source subsets.

\begin{table}[t]
\centering
\caption{Per-Database Performance on Deduplicated Test Split}
\label{tab:subset_dedup}
\begin{tabular}{lcccc}
\toprule
\textbf{Database} & \textbf{Samples} & \textbf{Accuracy (\%)} & \textbf{Normal Acc. (\%)} & \textbf{Abnormal Acc. (\%)} \\
\midrule
training-a &  61 & \pending & \pending & \pending \\
training-b &  74 & \pending & \pending & \pending \\
training-c &   5 & \pending & \pending & \pending \\
training-d &   8 & \pending & \pending & \pending \\
training-e & 321 & \pending & \pending & \pending \\
training-f &  17 & \pending & \pending & \pending \\
\midrule
\textbf{Overall} & \textbf{486} & \pending & \pending & \pending \\
\bottomrule
\end{tabular}
\end{table}


Recording site performance breakdowns highlight disparities between clinical recording environments.

%=====================================================================
\section{Limitations}
Several limitations qualify this study. First, evaluation is restricted to the PhysioNet/CinC 2016 corpus, which lacks multi-center longitudinal patient tracking. Second, the 77:23 class imbalance remains uncorrected during dataset construction and is managed solely via loss reweighting and balanced metrics. Third, the ResNet50 backbone was pretrained on natural images from ImageNet rather than domain-specific biomedical acoustic spectrograms. Fourth, PhysioNet state annotations are generated via semi-automated logistic regression hidden semi-Markov models \cite{springer} rather than manual expert delineations, introducing potential label noise into enrichment calculations. Fifth, the dilated backbone expands the spatial token sequence from 49 to 196, quadrupling self-attention computational complexity.

%=====================================================================
\section{Conclusion and Future Work}
This work examined phonocardiogram classification through five structured contributions: measuring data leakage from duplicate validation entries, evaluating architectural component ablations across three random seeds, implementing a dilated backbone narrowing temporal resolution to 0.57~s per cell, conducting a population-level explanation audit across five quantitative dimensions, and evaluating model calibration with sensitivity-targeted clinical threshold selection.

Future PCG research should adopt deduplicated data partitioning by default, report leave-one-database-out cross-site performance alongside in-distribution metrics, and conduct population-level faithfulness and sanity audits before relying on attribution visualizations. Future technical work will explore self-supervised acoustic pretraining on large audio repositories, joint end-to-end segmentation and classification architectures, and prospective clinical validation using calibrated digital stethoscopes.

%=====================================================================
\section*{Code Availability}
The implementation codebase, preprocessing scripts, model checkpoints, audit suites and evaluation scripts are available at: \url{https://github.com/urvagandhi/heart-sound-analysis}.

\section*{Acknowledgment}
The authors express their sincere gratitude to Dr.\ Sapan Mankad for guidance, insightful discussions and project supervision throughout this work.

%=====================================================================
\begin{thebibliography}{99}
\bibitem{who_cvd} World Health Organization, ``Cardiovascular diseases (CVDs),'' fact sheet. [Online]. Available: \url{https://www.who.int/news-room/fact-sheets/detail/cardiovascular-diseases-(cvds)}
\bibitem{rubin} J. Rubin, R. Abreu, A. Ganguli, S. Nelaturi, I. Matei and K. Sricharan, ``Classifying heart sound recordings using deep convolutional neural networks and mel-frequency cepstral coefficients,'' in \textit{Proc. Computing in Cardiology (CinC)}, vol. 43, 2016, pp. 813--816.
\bibitem{thomae} C. Thomae and A. Dominik, ``Using deep gated RNN with a convolutional front end for end-to-end classification of heart sound,'' in \textit{Proc. Computing in Cardiology (CinC)}, vol. 43, 2016, pp. 625--628.
\bibitem{springer} D. B. Springer, L. Tarassenko and G. D. Clifford, ``Logistic regression-HSMM-based heart sound segmentation,'' \textit{IEEE Trans. Biomed. Eng.}, vol. 63, no. 4, pp. 822--832, 2016.
\bibitem{alrabie} S. Alrabie and A. Barnawi, ``Are artificial intelligence models listening like cardiologists? Bridging the gap between artificial intelligence and clinical reasoning in heart-sound classification using explainable artificial intelligence,'' \textit{Bioengineering}, vol. 12, no. 6, p. 558, 2025.
\bibitem{adebayo} J. Adebayo, J. Gilmer, M. Muelly, I. Goodfellow, M. Hardt and B. Kim, ``Sanity checks for saliency maps,'' in \textit{Proc. NeurIPS}, vol. 31, 2018.
\bibitem{chakrabarty} S. Chakrabarty, P. Bishwas, M. Bandyopadhyay and J. Sublime, ``Can we trust AI with our ears? A cross-domain comparative analysis of explainability in audio intelligence,'' \textit{IEEE Access}, vol. 13, pp. 36221--36245, 2025.
\bibitem{he} K. He, X. Zhang, S. Ren and J. Sun, ``Deep residual learning for image recognition,'' in \textit{Proc. IEEE CVPR}, 2016, pp. 770--778.
\bibitem{hu} J. Hu, L. Shen and G. Sun, ``Squeeze-and-excitation networks,'' in \textit{Proc. IEEE CVPR}, 2018, pp. 7132--7141.
\bibitem{vaswani} A. Vaswani \textit{et al.}, ``Attention is all you need,'' in \textit{Proc. NeurIPS}, vol. 30, 2017.
\bibitem{selvaraju} R. R. Selvaraju \textit{et al.}, ``Grad-CAM: Visual explanations from deep networks via gradient-based localization,'' in \textit{Proc. IEEE ICCV}, 2017, pp. 618--626.
\bibitem{lundberg} S. M. Lundberg and S.-I. Lee, ``A unified approach to interpreting model predictions,'' in \textit{Proc. NeurIPS}, vol. 30, 2017.
\bibitem{li} F. Li, H. Tang, S. Shang, K. Mathiak and F. Cong, ``Classification of heart sounds using convolutional neural network,'' \textit{Appl. Sci.}, vol. 10, no. 11, Art. no. 3956, 2020.
\bibitem{ren} Z. Ren \textit{et al.}, ``Deep attention-based neural networks for explainable heart sound classification,'' \textit{Mach. Learn. Appl.}, vol. 9, Art. no. 100322, 2022.
\bibitem{padhy} S. K. Padhy, A. Mohapatra and S. Patra, ``X-CBNet: An explainable effective deep learning framework based on spectrograms for predicting valvular disorder using PCG signals,'' \textit{J. Transform. Technol. Sustain. Dev.}, vol. 9, Art. no. 18, 2025.
\bibitem{althaph} B. Althaph and N. P. Challa, ``Explainable attention-based deep learning for classification and interpretation of heart murmurs using phonocardiograms,'' \textit{Sci. Rep.}, vol. 15, Art. no. 37991, 2025.
\bibitem{ramakrishna} J. S. Ramakrishna, S. C. Venkateswarlu, K. N. Kumar and P. Shreya, ``Development of explainable machine intelligence models for heart sound abnormality detection,'' \textit{Indones. J. Electr. Eng. Comput. Sci.}, vol. 36, no. 2, pp. 846--853, 2024.
\bibitem{guleria} P. Guleria, P. N. Srinivasu, S. Ahmed, N. Almusallam and F. K. Alarfaj, ``XAI framework for cardiovascular disease prediction using classification techniques,'' \textit{Electronics}, vol. 11, no. 24, Art. no. 4086, 2022.
\bibitem{clifford} G. D. Clifford \textit{et al.}, ``Classification of normal/abnormal heart sound recordings: The PhysioNet/Computing in Cardiology Challenge 2016,'' in \textit{Proc. Computing in Cardiology}, vol. 43, 2016, pp. 609--612.
\bibitem{liu} C. Liu \textit{et al.}, ``An open access database for the evaluation of heart sound algorithms,'' \textit{Physiol. Meas.}, vol. 37, no. 12, pp. 2181--2213, 2016.
\bibitem{park} D. S. Park \textit{et al.}, ``SpecAugment: A simple data augmentation method for automatic speech recognition,'' in \textit{Proc. Interspeech}, 2019, pp. 2613--2617.
\bibitem{deng} J. Deng \textit{et al.}, ``ImageNet: A large-scale hierarchical image database,'' in \textit{Proc. IEEE CVPR}, 2009, pp. 248--255.
\bibitem{adamw} I. Loshchilov and F. Hutter, ``Decoupled weight decay regularization,'' in \textit{Proc. ICLR}, 2019.
\bibitem{sgdr} I. Loshchilov and F. Hutter, ``SGDR: Stochastic gradient descent with warm restarts,'' in \textit{Proc. ICLR}, 2017.
\bibitem{szegedy} C. Szegedy, V. Vanhoucke, S. Ioffe, J. Shlens and Z. Wojna, ``Rethinking the Inception architecture for computer vision,'' in \textit{Proc. IEEE CVPR}, 2016, pp. 2818--2826.
\bibitem{micikevicius} P. Micikevicius \textit{et al.}, ``Mixed precision training,'' in \textit{Proc. ICLR}, 2018.
\bibitem{guo} C. Guo, G. Pleiss, Y. Sun and K. Q. Weinberger, ``On calibration of modern neural networks,'' in \textit{Proc. ICML}, 2017, pp. 1321--1330.
\bibitem{yu} F. Yu and V. Koltun, ``Multi-scale context aggregation by dilated convolutions,'' in \textit{Proc. ICLR}, 2016.
\bibitem{petsiuk} V. Petsiuk, A. Das and K. Saenko, ``RISE: Randomized input sampling for explanation of black-box models,'' in \textit{Proc. BMVC}, 2018.
\end{thebibliography}

\end{document}
